%--------------------------------------------------------------------------------
%
%
%--------------------------------------------------------------------------------
\pagebreak
\section*{DIAG - Diagnose Instruction}
\addcontentsline{toc}{section}{DIAG - Diagnose Instruction}

%--------------------------------------------------------------------------------
\begin{iPageDescEntry}{Syntax}
\texttt{DIAG id, RegB, RegA}
\end{iPageDescEntry}

%--------------------------------------------------------------------------------
\begin{iPageDescEntry}{Format}
The \texttt{DIAG} instruction uses the signed immediate S9 instruction format. \\

\begin{flushleft}
\drawInstrFormat
    {0,1,3,5,6.5,8.5,9.5,11.5,16}
    {\OpCodeGrpSYS, \OpCodeDIAG, RegR, id1, RegB, id2, RegA, 0}
\end{flushleft}
\end{iPageDescEntry}

%--------------------------------------------------------------------------------
\begin{iPageDescEntry}{Description}
The \texttt{DIAG} instruction provides a mechanism for diagnostic and system level operations.  

The concatenated \texttt{id1} and \texttt{id2} fields specify the diagnostic function. The \texttt{RegB} and \texttt{RegA} fields may provide operands or data required by the diagnostic function. Any result is returned in \texttt{RegR}.
 
\end{iPageDescEntry}

%--------------------------------------------------------------------------------
\begin{iPageDescOpEntry}{Operation}
\begin{lstlisting}[style=iPageOpStyle]
RegR <- diagOperation( id, RegB, RegA );
\end{lstlisting}
\end{iPageDescOpEntry}

%--------------------------------------------------------------------------------
\begin{iPageDescEntry}{Exceptions}
Exceptions depend on diagnostic function and operands.
\end{iPageDescEntry}

%--------------------------------------------------------------------------------
\begin{iPageDescEntry}{Notes}
None
\end{iPageDescEntry}
